\subsection{Linear Combinations: What Addition and Scaling Can Build}

Addition and scalar multiplication were introduced as two separate operations,
but they are most useful when they are combined. Given vectors $\mathbf u$ and
$\mathbf v$ and real numbers $a$ and $b$, the vector
\[
a\mathbf u+b\mathbf v
\]
is called a \emph{linear combination} of $\mathbf u$ and $\mathbf v$.

The phrase has a direct geometric meaning. First scale $\mathbf u$ by $a$ and
$\mathbf v$ by $b$; then add the two resulting displacements. Thus a linear
combination is not a new operation. It is a controlled use of the two operations
already understood.

\begin{center}
\begin{tikzpicture}[scale=0.9,>=stealth]
    \coordinate (O) at (0,0);
    \coordinate (AU) at (3.4,0.9);
    \coordinate (BV) at (1.0,2.7);
    \coordinate (S) at (4.4,3.6);

    \draw[blue!70!black,line width=1.1pt,->] (O) -- (AU)
        node[midway,below] {$a\mathbf u$};
    \draw[orange!85!black,line width=1.1pt,->] (AU) -- (S)
        node[midway,right] {$b\mathbf v$};
    \draw[very thick,->] (O) -- (S)
        node[midway,above left] {$a\mathbf u+b\mathbf v$};
    \draw[helper,dashed] (O) -- (BV);
    \draw[helper,dashed] (BV) -- (S);
\end{tikzpicture}
\end{center}

A single nonzero vector can generate only one line through the origin:
\[
\{\lambda\mathbf v:\lambda\in\mathbb R\}.
\]
Changing $\lambda$ changes the length and possibly the orientation, but never
the underlying line of motion.

Two vectors behave differently. If $\mathbf u$ and $\mathbf v$ are parallel,
then every linear combination still lies on one line. If they are not parallel,
then their linear combinations fill the whole plane. Geometrically, the two
independent directions allow us to reach any endpoint by a suitable
parallelogram.

\begin{theorembox}
Two nonparallel vectors $\mathbf u$ and $\mathbf v$ determine every vector in
the plane uniquely in the form
\[
\mathbf w=a\mathbf u+b\mathbf v.
\]
The numbers $a$ and $b$ record how much motion occurs in each chosen direction.
\end{theorembox}

This observation is the bridge to the later notion of a basis. A basis is not
an unrelated new object: it is a pair of directions whose linear combinations
can describe every vector uniquely.